\chapter{Introduction} \label{ch:intro}

Remote sensing performed with objective of surface observation typically has to compensate for the optical effects of the intervening atmosphere between the surface and the sensor. The exceptions are qualitative retrievals within a scene, or quantitative retrievals over targets that are bright at a given spectral window and at geometries and spectral windows where the atmosphere has a low optical thickness. The optical effects of the atmosphere can be quantified in terms of transmission and reflectance functions for the atmospheric layer, arising out of absorption and scattering by gases and aerosols in the path between the light source (the Sun for passive remote sensing), the surface, and the sensor.

Several processors exist for compensation of atmospheric effects on remote imagery, though most processors compensate only for the magnitude of the effects, without explicitly compensating for the blurring effect caused by multiple interaction of light between the surface and the atmosphere and the upward diffuse transmittance of the atmosphere for the path between the surface and the sensor. This blurring effect, also know as adjacency effect, is a function of the reflectance contrast, size and spatial distribution of surface types and the magnitude and spatial scale of the atmospheric scattering contributing to the recorded signal. In general, the higher the contrast between a target and its surroundings, the smaller the target area, and the larger the atmospheric turbidity, the larger the blurring effect of the atmosphere will be.

The spatial scale of the effect depends on the major atmospheric scatter. Gases have nearly isotropic scattering phase functions (Rayleigh scattering), resulting that light reflected at high zenith angles from surrounding surfaces (and therefore from further away from the target), has a similar probability of being scattered towards the sensor as from light reflected at low zenith angles. In contrast, the scattering phase function of aerosols has a forward peak, resulting in lower probability that light reflected at high zenith angles from surrounding surfaces will be scattered towards the sensor, when compared to light reflected at low zenith angles. The spatial scale of the effect can be quantified as the radius from which 100~\% of the diffusely transmitted signal arises. For a nadir looking sensor at the top of a Rayleigh atmosphere this radius is in the order of 50~km, while for a purely (hypothetical) aerosol atmosphere this radius is in the order of 10~km, with higher contributions from surfaces closer to the target. The relative dominance of each scatter type for the blurring effect of the atmosphere is given by the ratio of the Rayleigh and aerosol contributions to the diffuse upward transmittance and spherical albedo, which in turn will vary with the spectral window, the atmospheric pressure at the surface, the concentration and vertical profile of the aerosols, the physical properties of those aerosols, and the optical thickness of each component.

Therefore, over regions with high contrast in surface reflectance, the blurring effect will typically be significant in the ultraviolet and blue end of the spectrum (due to Rayleigh scattering) and at longer wavelengths depending on the time-variable aerosol composition and concentration. Examples of surfaces with high contrast are the transitions between water bodies and land at visible wavelengths longer than the red. Since size is also a significant component of the blurring effect, and most rivers have widths $< 1$~km \citep{Dekker2018} and most lakes have areas $< 1$~km$^2$ \citep{Verpoorter2014}, the blurring effect of the atmosphere can also be significant even in the visible range for the majority of the inland water bodies. Correction of this effect is therefore essential for the accurate retrieval of water quality parameters (\textit{e.g.} water turbidity and pigment concentration) in such water bodies, both for the exploitation of historical archives and future acquisitions. Compensation for the blurring effect, also known as adjacency correction, is hence crucial for long-term analysis of satellite remote sensing data, especially of inland water quality. Such processing was considered a key research development for the aquatic remote sensing community (Water ForCE 2023: Water Quality Continuum Atmospheric Correction Recommendations Questionnaire Results; \url{https://waterforce.eu/}).

The main objective of RAdCor is to provide a physically-based, sensor- and scene-generic, computationally efficient, free an open source processor to compensate for atmospheric effects in remote optical imagery that accounts for the blurring effect. This document describes the approach developed in the RAdCor project, its theoretical background, and specific implementation in the free and open source software ACOLITE \citep[][\url{https://odnature.naturalsciences.be/remsem/software-and-data/acolite}]{Vanhellemont2018,Vanhellemont2019a,Vanhellemont2019b}. Strictly, the analysis presented in Chapters~\ref{ch:sim01} and \ref{ch:sim02} cover a version prior to the public version release. A separate document was submitted to peer-review publication containing the main topics of this document and validation against imagery and in situ data \citep{Castagna2024} using the first public release version.

